% !PRINTED 522
\weberpage{522}

Um aber eine Anwendung unserer allgemeinen Methode zu geben, nehmen wir $m \equiv 3\ (\mathrm{mod}\ 8)$ an und setzen $D = -m$. Wir befriedigen die Gleichung:
\[4n = y^2 + mx^2,\]
indem wir setzen
\[m = m'm'',\]
\[x = 1, \quad y = \frac{m'-m''}{2}, \quad 4n = \left(\frac{m'+m''}{2}\right)^2,\]
wobei $y$ und folglich $n$ ungerade ausfällt. Es wird ferner nach § 136, (13)
\[\alpha + \beta\omega = \left(\frac{\sqrt{m'} + i\sqrt{m''}}{2}\right)^2.\]

Setzen wir unter der Voraussetzung, daß $n$ durch 3 nicht teilbar ist,
\weq{1}{M = \left(\frac{c}{e}\right) i^{\frac{a-1}{2}} \sqrt{\partial}\, \frac{\eta\left(\dfrac{c+\partial\omega}{a}\right)}{\eta(\omega)},}
so läßt sich $M$ mit Anwendung von § 38, (15) bestimmen und man findet, wenn $m''$ denjenigen der beiden Faktoren $m'$, $m''$ bedeutet, der modulo 4 mit 1 kongruent ist:
\weq{2}{M = -\left(\frac{A}{m''}\right)(-1)^{\frac{m'-m''-2}{8}} e^{-\frac{\pi i}{3}A(m'-m'')}\, \frac{\sqrt{m''} - i\sqrt{m'}}{2}.}
Um hiervon eine Anwendung zu machen, setzen wir $n = 25$, also
\[20 = m' + m'',\]
und folglich
\[
\begin{aligned}
m'' &= 1,  & m' &= 19, & m &= 19,\\
m'' &= 17, & m' &= 3,  & m &= 51,\\
m'' &= 13, & m' &= 7,  & m &= 91,\\
m'' &= 9,  & m' &= 11, & m &= 99.
\end{aligned}
\]
Legen wir die Form
\[\left(1, \ \frac{m'-m''}{2}, \ n\right)\]
zugrunde, so ist $A = 1$ zu setzen, und (2) ergibt:
\weq{3}{\begin{array}{ll}
m = 19, & M = -\dfrac{1-i\sqrt{19}}{2},\\[6pt]
m = 51, & M = e^{-\frac{\pi i}{3}}\dfrac{\sqrt{17} - i\sqrt{3}}{2},\\[6pt]
m = 91, & M = \dfrac{\sqrt{13} - i\sqrt{7}}{2},\\[6pt]
m = 99, & M = e^{\frac{\pi i}{3}}\dfrac{3 - i\sqrt{11}}{2}.
\end{array}}
